\begin{figure*}[t]
    \centering
    \includegraphics[width=\linewidth]{figures/fig2_method.pdf}
    \caption{
        (a) A unified view comparing LAP-3B with prior VLAs in terms of action representation. A VLM backbone predicts discrete language-action tokens using a cross-entropy objective, while a lightweight action expert predicts continuous actions via flow matching. Gradients from the action expert are blocked from the VLM through knowledge insulation \cite{driess2025knowledgeinsulatingvisionlanguageactionmodels}, ensuring that the VLM is trained purely via language supervision. At test time, the action expert is rolled out for fast inference.
        (b) Visualizations of language-actions from the DROID dataset \cite{khazatsky2025droidlargescaleinthewildrobot}.
        }
    \label{fig:method}
\end{figure*}


\section{LAP: Language-Action Pre-Training for VLAs}
\label{sec:method}

We now introduce \emph{Language-Action Pre-training} (LAP), a general pre-training recipe for VLAs. We first formalize LAP's language-action learning objective (Section~\ref{subsec:problem_formulation}), then describe how to obtain language-actions at scale without manual annotation (Section~\ref{subsec:curation}). We then present \textsc{LAP-3B}, a vision-language-action model trained with LAP that demonstrates substantial cross-embodiment transfer (Section~\ref{subsec:architecture}), followed by specific implementation details (Section~\ref{sec:implementation_details}).

\subsection{Formulations}
\label{subsec:problem_formulation}

LAP reframes action learning in VLAs as a standard conditional language modeling problem. Given an observation $o_t = \{I_t^1, \dots, I_t^n, s_t\}$ consisting of multi-view RGB images and proprioceptive state, and a task instruction $l$, we deterministically derive a structured language-action sequence $\hat{a}_{t:t+H}^{\text{lang.}}$, as shown in Fig.\ref{fig:method}$(b)$, summarizing the corresponding continuous action chunk $a_{t:t+H}$.

The VLM backbone is trained to auto-regressively predict the language-action tokens conditioned on $(o_t, l)$ using the standard cross-entropy objective
\begin{equation}
\mathcal{L}_{\text{CE.}}
=
-\mathbb{E}_{(\hat{a}_{t:t+H}^{\text{lang.}}, o_t, l) \sim \mathcal{D}}
\sum_{i}
\log p_{\theta}
\big(
\hat{a}^{\text{lang.}}_{t,i}
\mid o_t, l, \hat{a}^{\text{lang.}}_{t,<i}
\big),
\label{eq:language-action}
\end{equation}
which is identical to standard sequence-to-sequence language modeling and corresponds to maximum likelihood estimation of language-actions, analogous to behavioral cloning for continuous actions.

\subsection{Language-Action Dataset Curation}
\label{subsec:curation}

A key requirement for LAP is the access to language-action supervision at scale. Rather than relying on manual annotation, we leverage the structured nature of robotic datasets to derive language-actions directly from end-effector trajectories.

Specifically, we represent language-actions as natural language descriptions of end-effector delta motions. Under a fixed coordinate convention and a small set of structured templates, continuous actions can be deterministically converted into language descriptions with arbitrary numerical resolution.

\noindent\textbf{Coordinate System.}
We adopt a convention commonly used in robotics datasets, where $+x$ is forward, $+y$ is left, and $+z$ is up; rotational components are expressed as Euler angles following the right-hand rule. Since most manipulation datasets already follow this convention, language-actions can be parsed directly without additional frame transformations. 

To encourage embodiment-agnostic reasoning and consistent use of multi-view visual observations, we additionally express language-actions in multiple reference frames. Specifically, during training we randomize the reference frame used to generate each language-action, expressing 50\% in the robot base frame and 50\% in the end-effector frame, and include the name of the reference frame in the model prompt. 

\noindent\textbf{Templates.}
Language-actions are generated using a small set of structured natural language templates that describe translational and rotational end-effector motions:
\begin{equation}
\hat{a}_t^{\text{lang.}}
=
\text{``}\langle\text{verb}\rangle\,
\langle\text{direction}\rangle\,
\langle\text{magnitude}\rangle\,
\langle\text{unit}\rangle\text{''},
\end{equation}
where the verb determines the action category: \emph{move} denotes translational end-effector displacement, \emph{tilt} denotes roll and pitch rotations, and \emph{rotate} denotes yaw rotation. For translation, directions correspond to $\{\pm x, \pm y, \pm z\}$ (e.g., forward/backward, left/right, up/down). For roll and pitch, we use ``left/right'' and ``back/forward'' following the right-hand rule convention. For yaw, ``clockwise'' and ``counterclockwise'' represent negative and positive rotations, respectively.


Rather than generating language-actions at every timestep, we represent $\hat{a}_{t:t+H}^{\text{lang.}}$ as a single description summarizing the net displacement of the action chunk $a_{t:t+H}$. Since actions are expressed as delta end-effector poses, this corresponds to the cumulative translation and rotation between timesteps $t$ and $t+H$. This temporal abstraction produces a low-frequency supervision signal that is easier for VLMs to model while abstracting away high-frequency control. We also discretize magnitudes to integer values and express translations in centimeters (e.g., \emph{``move forward 5 cm''}), yielding language-actions that are both structured and precise. Visual examples are shown in Fig.~\ref{fig:method}(b).

Compared to prior action tokenization approaches that discretize continuous actions into arbitrary or out-of-vocabulary symbols, language-actions remain well supported by the VLM’s natural language pre-training distribution, mitigating the distributional mismatch known to degrade representation quality \cite{hancock2025actionslanguagefinetuningvlms, driess2025knowledgeinsulatingvisionlanguageactionmodels}. In contrast to learned tokenizers such as FAST \cite{pertsch2025fastefficientactiontokenization}, which introduce a separate representation learning stage to compress actions into discrete non-linguistic tokens, language-actions are produced purely through deterministic parsing, eliminating the need for any learned action tokenizer. Moreover, unlike conventional action representations, language-actions naturally align with the input–output format of Visual--Question--Answering (VQA) tasks, enabling synergistic co-training. For instance, in a ``motion prediction'' VQA task, the model is prompted with image pairs $(I_t^i, I_{t+H}^i)$ and trained to predict the corresponding language-action $\hat{a}_{t:t+H}^{\text{lang.}}$ describing the observed displacement.


\subsection{Model Architecture and Training}
\label{subsec:architecture}

LAP is agnostic to the underlying VLA architecture and can, in principle, be integrated into any VLM-based policy by treating action prediction as a language modeling problem. However, auto-regressive generation of language-actions at inference time is slow and impractical for real-time robot control. To obtain an efficient system, we instantiate LAP as LAP-3B, which adopts a Mixture-of-Transformers architecture \cite{liang2024mixture} combining a LAP-trained VLM backbone with a lightweight flow-matching action expert for high-frequency motor execution, following $\pi_{0.5}$ \cite{intelligence2025pi05visionlanguageactionmodelopenworld}. This design enables a controlled comparison: \textsc{LAP-3B} differs from $\pi_{0.5}$ only in the action representation used to supervise the VLM (language-actions versus FAST tokens), isolating the downstream effect of the representation itself, which we evaluate in Sec.~\ref{sec:experiments}.

During training, the VLM backbone is optimized to predict structured language-actions using the cross-entropy loss in Eq.~\eqref{eq:language-action}, while the action expert predicts continuous action chunks $a_{t:t+H}$ via a flow-matching objective
\begin{equation}
\mathcal{L}_{\text{FM}}
=
\mathbb{E}_{(a_{t:t+H}, o_t, l) \sim \mathcal{D}}
\left[
\left\|
v_{\phi}(o_t, l, a_{t:t+H}, \tau) - \dot{x}_{\tau}
\right\|^2
\right],
\end{equation}
where $v_{\phi}$ denotes the velocity field parameterized by the action expert and $\dot{x}_{\tau}$ is the target flow corresponding to the ground-truth action chunk (see Appendix~\ref{app:implementation_details} for the full formulation). The overall training objective is
\begin{equation}
\mathcal{L}
=
\mathcal{L}_{\text{FM}}
+
\lambda\,\mathcal{L}_{\text{CE}}.
\end{equation}

Under this architecture, the VLM and action expert communicate solely through cross-attention. We apply causal masking to ensure that the raw action and the language-action do not attend to each other.
Following \cite{driess2025knowledgeinsulatingvisionlanguageactionmodels}, we further block gradients from the action expert propagating back into the VLM backbone. While not required by LAP, this knowledge insulation helps preserve pre-trained VLM representations and stabilizes joint training. At inference time, only the action expert is rolled out for control, enabling real-time execution at 25\,Hz on an NVIDIA RTX~4090 GPU.



\subsection{Implementation Details}
\label{sec:implementation_details}


\textsc{LAP-3B}'s VLM backbone is initialized as the PaliGemma-3B \cite{beyer2024paligemmaversatile3bvlm} model. We train on the Open X-Embodiment (OXE) dataset \cite{embodimentcollaboration2025openxembodimentroboticlearning} and MolmoAct dataset \cite{lee2025molmoactactionreasoningmodels} (see full data mixture details in Appendix~\ref{app:implementation_details}). We use a shuffle buffer size of 16 million samples to ensure near uniform mixing, which improves training stability.
For proprioception, we represent state $s_t$ using Cartesian end-effector pose: position, a continuous 6D rotation representation \cite{zhou2019continuity}, and a binary gripper state. Since open-source datasets use inconsistent rotation conventions, we canonicalize rotations to extrinsic XYZ before converting to the 6D representation. The proprioceptive state is discretized into 255 bins and appended to the prompt as text tokens, following \cite{intelligence2025pi05visionlanguageactionmodelopenworld}. Actions are represented as delta end-effector poses. We normalize states and actions using the 1st and 99th global quantiles computed over the full training dataset, and reuse the same statistics at inference on unseen embodiments.

Training uses a global batch size of 2048 across 64 TPU v6e chips for 15k gradient steps (\emph{approximately only 10 wall-clock hours}), corresponding to roughly 0.65 epoch over the full dataset, at which point the checkpoint already performs well in real-world experiments. Exponential moving average (EMA) updates begin after 5k steps.
Image inputs are resized to $224 \times 224$, with at most two images per sample. We use a fixed learning rate of $1 \times 10^{-4}$ with linear warmup over the first 5,000 steps.

